\section{Related Work}
\label{sec:related}

\subsection{EEG Artifact Correction}
Classical ocular-artifact correction includes regression, filtering, and blind
source separation.  ICA is effective when artifact components are separable,
but component selection and decomposition stability can vary with montage and
recording duration.  Regression methods use EOG references more directly, yet
their transfer coefficients can differ strongly across participants and
sessions.  Deep denoisers trained on paired or semi-simulated examples avoid
manual component labeling, including the EEGdenoiseNet line of work
\cite{zhang2021eegdenoisenet}, but usually apply one population mapping at
deployment.

Our work retains the interpretable EOG-to-EEG transfer while using a generative
model only for low-dimensional artifact coefficients.  The support recording is
query-disjoint, and query EOG is reserved for evaluation.  This differs from
methods that require an EOG channel throughout deployment.

\subsection{Cross-Participant Adaptation}
EEG varies across participants because of anatomy, impedance, reference,
behavior, and acquisition state.  Many subject-aware networks encode a training
participant label or learn a closed-set identity embedding.  Such an embedding
cannot directly represent a previously unseen user and risks personalizing
neural content rather than the acquisition corruption.  Domain-adaptation
methods address population shift but typically do not identify the causal role
of a short support calibration.

We instead personalize only the pollution operator.  A target participant is
represented by a support-derived artifact basis, not by an ID.  Matching,
population, and same-cell wrong interventions share network weights and queries,
which makes the subject-calibration increment directly testable.

\subsection{Diffusion for Biosignal Denoising}
Diffusion models offer stable denoising objectives and posterior samples without
an adversarial discriminator \cite{ho2020ddpm}.  Existing biosignal applications
include ECG reconstruction \cite{zhao2023ecg} and noisy-signal-conditioned EEG
models.  Conditioning may describe the noisy observation, an artifact type, a
domain representation, or a participant label.  These alternatives do not by
themselves ensure that a model uses an unseen user's calibration; a sufficiently
strong network may ignore a soft context vector.

The present construction makes support information structural.  The calibrated
basis defines the diffusion coordinate system, the observation condition at
every reverse step, and the reconstruction operator.  It also differs from
full-EEG diffusion: only two bounded artifact-coefficient channels are diffused,
and the observed orthogonal complement is copied exactly.  Finally, we compare
against a capacity-matched one-step model trained on the identical coefficient
target, preventing a full-generation-versus-regression mismatch from being
misinterpreted as diffusion utility.
